\section*{第 \ref{ch:g12:complex} 章——复数}

\begin{solution}{exo:g12:complex:1}
$(3-2\iu)(1+4\iu) = 3 + 12\iu - 2\iu - 8\iu^2 = 11 + 10\iu$.

$\dfrac{2+\iu}{1-\iu} = \dfrac{(2+\iu)(1+\iu)}{(1-\iu)(1+\iu)}
= \dfrac{1 + 3\iu}{2} = \dfrac12 + \dfrac32\iu$.

$\iu^{2027}$：由 $\iu^4 = 1$ 及 $2027 = 4 \times 506 + 3$，
$\iu^{2027} = \iu^3 = -\iu$.
\end{solution}

\begin{solution}{exo:g12:complex:2}
\emph{(a)} $\Delta = 16 - 52 = -36$：
$z = \dfrac{4 \pm 6\iu}{2} = 2 \pm 3\iu$。二者为\omterm{def:g12:complex:conjugate}{共轭复数}，其积
$= 4 + 9 = 13 = \frac ca$。\emph{(b)} $\Delta = 1 - 4 = -3$：
$z = \dfrac{-1 \pm \iu\sqrt3}{2}$（即单位立方根 $j$ 与
$\conj j$），其积 $= \frac{1 + 3}{4} = 1$。
\end{solution}

\begin{solution}{exo:g12:complex:3}
$\abs{z_1} = \sqrt2$，\omterm{def:g12:complex:argument}{辐角}为 $\frac{3\pi}{4}$（第二象限）：
$z_1 = \sqrt2\,\eu^{3\iu\pi/4}$.
$\abs{z_2} = 2$，\omterm{def:g12:complex:argument}{辐角}为 $-\frac\pi6$：
$z_2 = 2\,\eu^{-\iu\pi/6}$。于是
\[
\frac{z_1}{z_2} = \frac{\sqrt2}{2}\,\eu^{\iu\left(\frac{3\pi}{4} + \frac{\pi}{6}\right)}
= \frac{\sqrt2}{2}\,\eu^{11\iu\pi/12}.
\]
用\omterm{def:g12:complex:def}{代数形式}计算，
\[
\frac{z_1}{z_2} = \frac{-1+\iu}{\sqrt3 - \iu}
= \frac{(-1+\iu)(\sqrt3+\iu)}{4}
= \frac{-\sqrt3 - 1 + \iu(\sqrt3 - 1)}{4}.
\]
比较\omterm{def:g12:complex:def}{实部}可得：
$\frac{\sqrt2}{2}\cos\frac{11\pi}{12} = \frac{-\sqrt3-1}{4}$，于是
\[
\cos\frac{11\pi}{12} = -\frac{\sqrt6 + \sqrt2}{4}.
\]
\end{solution}

\begin{solution}{exo:g12:complex:4}
\emph{(a)} $\abs{z - 2} = \abs{z - (-\iu)}$ 表示点 $M$ 到
$A(2, 0)$ 与 $B(0, -1)$ 等距，即 $[AB]$ 的\emph{中垂线}。

\emph{(b)} 距离为 $3$ 的点（参照点为 $(1,1)$）构成\emph{圆}，
圆心为 $1 + \iu$，半径为 $3$。

\emph{(c)} 由 \cref{prop:g12:complex:geometry}(4)，$\frac{z-1}{z+1}$ 为纯
虚数，即从 $A(1,0)$ 到 $M$ 的\omterm{def:g11:vect:vector}{向量}与从 $B(-1,0)$ 到 $M$ 的\omterm{def:g11:vect:vector}{向量}
\omterm{def:g11:scal:orthogonal}{正交}，即 $M$ 对线段 $[AB]$ 张直角。所求集合为
\emph{以 $[AB]$ 为直径的圆}（即单位圆）\emph{去掉
点 $A$ 与 $B$ 本身}。
\end{solution}

\begin{solution}{exo:g12:complex:5}
棣莫弗公式：$\cos3\theta + \iu\sin3\theta = (\cos\theta + \iu\sin\theta)^3$。
由二项式定理\omterm{def:g10:algebra:expand}{展开}，记 $c = \cos\theta$，$s = \sin\theta$：
\[
(c + \iu s)^3 = c^3 + 3\iu c^2 s - 3 c s^2 - \iu s^3
= (c^3 - 3cs^2) + \iu(3c^2 s - s^3).
\]
比较\omterm{def:g12:complex:def}{实部}与\omterm{def:g12:complex:def}{虚部}，并利用 $s^2 = 1 - c^2$，$c^2 = 1 - s^2$：
\[
\cos3\theta = 4\cos^3\theta - 3\cos\theta,
\qquad
\sin3\theta = 3\sin\theta - 4\sin^3\theta .
\]
\end{solution}

\begin{solution}{exo:g12:complex:6}
\[
\cos^3\theta = \left(\frac{\eu^{\iu\theta} + \eu^{-\iu\theta}}{2}\right)^{\!3}
= \frac{\eu^{3\iu\theta} + 3\eu^{\iu\theta} + 3\eu^{-\iu\theta} + \eu^{-3\iu\theta}}{8}
= \frac{\cos3\theta + 3\cos\theta}{4}.
\]
由此
\[
\int_0^{\pi/2}\cos^3\theta\,\dd\theta
= \frac14\left[\frac{\sin3\theta}{3} + 3\sin\theta\right]_0^{\pi/2}
= \frac14\left(-\frac13 + 3\right) = \frac{2}{3}.
\]
（与 \cref{exo:g12:integ:9} 中的 $W_3 = \frac23$ 一致，因为
$\int_0^{\pi/2}\cos^3 = \int_0^{\pi/2}\sin^3$，由对称性
$\theta \mapsto \frac\pi2 - \theta$ 得到）。
\end{solution}

\begin{solution}{exo:g12:complex:7}
$c = a + \eu^{\pm\iu\pi/3}(b - a)$，其中 $b - a = 2 + \iu$，且
$\eu^{\pm\iu\pi/3} = \frac12 \pm \iu\frac{\sqrt3}{2}$：
\[
\eu^{\iu\pi/3}(2+\iu) = \left(\tfrac12 + \iu\tfrac{\sqrt3}{2}\right)(2+\iu)
= \left(1 - \tfrac{\sqrt3}{2}\right) + \iu\left(\tfrac12 + \sqrt3\right),
\]
于是 $c_1 = \left(2 - \frac{\sqrt3}{2}\right) +
\iu\left(\frac32 + \sqrt3\right)$，类似地，取\omterm{def:g12:complex:argument}{辐角}
$-\frac\pi3$：
$c_2 = \left(2 + \frac{\sqrt3}{2}\right) + \iu\left(\frac32 - \sqrt3\right)$.
\end{solution}

\begin{solution}{exo:g12:complex:8}
$-4 = 4\,\eu^{\iu\pi}$，故\omterm{def:g10:algebra:equation}{方程} $z^4 = -4$ 的解为
\[
z_k = \sqrt2\,\eu^{\iu\left(\frac\pi4 + \frac{k\pi}{2}\right)},
\quad k = 0, 1, 2, 3,
\]
\emph{即} $1 + \iu$、$-1 + \iu$、$-1 - \iu$、$1 - \iu$，它们是
外接圆半径为 $\sqrt 2$ 的正方形的顶点。将共轭根两两分组，
\[
z^4 + 4 = \bigl(z - (1{+}\iu)\bigr)\bigl(z - (1{-}\iu)\bigr)
\bigl(z + (1{-}\iu)\bigr)\bigl(z + (1{+}\iu)\bigr)
= (z^2 - 2z + 2)(z^2 + 2z + 2).
\]
\end{solution}

\begin{solution}{exo:g12:complex:9}
设 $q = \eu^{\iu\theta} \neq 1$：
\[
\sum_{k=0}^{n} \eu^{\iu k\theta} = \frac{\eu^{\iu(n+1)\theta} - 1}{\eu^{\iu\theta} - 1}
= \frac{\eu^{\iu(n+1)\theta/2}}{\eu^{\iu\theta/2}}
\cdot\frac{\eu^{\iu(n+1)\theta/2} - \eu^{-\iu(n+1)\theta/2}}{\eu^{\iu\theta/2} - \eu^{-\iu\theta/2}}
= \eu^{\iu n\theta/2}\,
\frac{\sin\frac{(n+1)\theta}{2}}{\sin\frac{\theta}{2}},
\]
利用 $\eu^{\iu\alpha} - \eu^{-\iu\alpha} = 2\iu\sin\alpha$（
\emph{半角分解}），取\omterm{def:g12:complex:def}{实部}：
\[
S = \frac{\sin\frac{(n+1)\theta}{2}}{\sin\frac{\theta}{2}}\,
\cos\frac{n\theta}{2}.
\]
\end{solution}

\begin{solution}{exo:g12:complex:10}
\emph{1.} $\omega$ 是 $5$ 次\omterm{def:g12:complex:rootsofunity}{单位根}，且异于 $1$，故
由\cref{thm:g12:complex:rootsofunity}，其五个根之和为零，即
$1 + \omega + \omega^2 + \omega^3 + \omega^4 = 0$.

\emph{2.} $u + v = \omega + \omega^2 + \omega^3 + \omega^4 = -1$，由第 1 问可得。
对于乘积，利用 $\omega^5 = 1$：
\[
uv = (\omega + \omega^4)(\omega^2 + \omega^3)
= \omega^3 + \omega^4 + \omega^6 + \omega^7
= \omega^3 + \omega^4 + \omega + \omega^2 = -1 .
\]
于是 $u, v$ 是\omterm{def:g10:algebra:equation}{方程} $X^2 + X - 1 = 0$ 的根：
$\left\{\frac{-1+\sqrt5}{2}, \frac{-1-\sqrt5}{2}\right\}$。注意到
$u = \omega + \conj\omega = 2\cos\frac{2\pi}{5} > 0$（因为
$\frac{2\pi}{5} < \frac{\pi}{2}$），故 $u = \frac{-1 + \sqrt5}{2}$。

\emph{3.} $\cos\frac{2\pi}{5} = \frac u2 = \frac{\sqrt5 - 1}{4}$.
\end{solution}

\begin{solution}{pb:g12:complex:1}
\textbf{1.} $(2 + \iu)(3 - \iu) = 6 - 2\iu + 3\iu + 1 =
7 + \iu$. \quad $\frac{1 + \iu}{1 - \iu} =
\frac{(1 + \iu)^2}{2} = \frac{2\iu}{2} = \iu$. \quad
$\abs{3 + 4\iu} = 5$.

\textbf{2.} $1 + \iu = \sqrt2\,\eu^{\iu\pi/4}$;
$-2 = 2\eu^{\iu\pi}$;
$1 - \iu\sqrt3 = 2\eu^{-\iu\pi/3}$.

\textbf{3.} $(1 + \iu)^8 = \left(\sqrt2\right)^8
\eu^{8\iu\pi/4} = 16\,\eu^{2\iu\pi} = 16$.

\textbf{4.} $z^2 = \eu^{\iu\pi/2}$:
$z = \pm\eu^{\iu\pi/4} = \pm\frac{\sqrt2}{2}(1 + \iu)$。又
$\Delta = 4 - 20 = -16$: $z = 1 \pm 2\iu$.

\textbf{5.} $z \mapsto \iu z$：旋转四分之一周，
绕原点。$\abs{z - (1 + \iu)} = 2$：圆心为
$1 + \iu$、半径为 $2$ 的圆。

\textbf{6.} $\left(\omega^k\right)^5 = \eu^{2\iu\pi k} = 1$,
且五个点 $\eu^{2\iu k\pi/5}$ 互不相同，故
这五个数恰为 $z^5 = 1$ 的全部解。它们位于单位圆上，
彼此以 $72^\circ$ 的等角相隔，构成一个正五边形，其一个顶点为
$1$.

\textbf{7.} $S = 1 + \omega + \dots + \omega^4$ 满足
$(1 - \omega)S = 1 - \omega^5 = 0$，且 $\omega \neq 1$，
故 $S = 0$。（正五边形的重心即其中心。）

\textbf{8.} $\omega^4 = \overline\omega$ 与
$\omega^3 = \overline{\omega^2}$：将共轭两两配对，
和的\omterm{def:g12:complex:def}{实部}即
$1 + 2\cos\frac{2\pi}{5} + 2\cos\frac{4\pi}{5} = 0$.

\textbf{9.} $\cos\frac{4\pi}{5} = 2c^2 - 1$，于是
$1 + 2c + 4c^2 - 2 = 0$，即 $4c^2 + 2c - 1 = 0$：
$c = \frac{-1 + \sqrt5}{4}$（取正根，因为
$72^\circ$ 为锐角）：
$\cos 72^\circ = \frac{\sqrt5 - 1}{4}$.

\textbf{10.} 数值上 $c \approx 0.30902$：与
计算器相符。又 $\frac{1}{2\varphi} = \frac{1}{1 + \sqrt5} =
\frac{\sqrt5 - 1}{4}$（分母有理化）：$c = \frac{1}{2\varphi}$。
正五边形处处体现黄金比：其对角线将
一边按比 $\varphi$ 分割——五角星的无穷
自相似性。

\textbf{11.} 平方根可用尺规作出（即
初中分册中几何平均的半圆作图法），
因此仅由\omterm{def:g10:numbers:sets}{有理数}经平方根运算构造的数
——如 $\cos 72^\circ$——是可作图的；立方根则
不然，而 $\cos 20^\circ$ 需要立方根：正五边形可作，
三等分角却做不到。高斯的判别准则决定了每一个正多边形
是否可作图——他 19 岁时作出的 17 边形
使他选择了数学。

\textbf{12.} $-4 = 4\eu^{\iu\pi}$：四次方根的
模为 $\sqrt2$，\omterm{def:g12:complex:argument}{辐角}为
$\frac{\pi}{4} + k\frac{\pi}{2}$：这四个数为
$\pm 1 \pm \iu$。将共轭两两配对：
$(z - 1 - \iu)(z - 1 + \iu) = z^2 - 2z + 2$ 与
$(z + 1 - \iu)(z + 1 + \iu) = z^2 + 2z + 2$：
即\cref{exo:g12:complex:8}中的\omterm{def:g10:algebra:expand}{因式分解}。

\textbf{13.} $(\cos\theta + \iu\sin\theta)^3 = \cos^3\theta +
3\iu\cos^2\theta\sin\theta - 3\cos\theta\sin^2\theta -
\iu\sin^3\theta$；由棣莫弗公式，此式等于
$\cos 3\theta + \iu \sin 3\theta$。取\omterm{def:g12:complex:def}{实部}：
$\cos 3\theta = \cos^3\theta - 3\cos\theta(1 - \cos^2\theta)
= 4\cos^3\theta - 3\cos\theta$——而\omterm{def:g12:complex:def}{虚部}则顺带给出
$\sin 3\theta$。

\textbf{14.} 它是
$1 + \omega + \dots + \omega^4 = 0$ 的\omterm{def:g12:complex:def}{实部}：和为 $0$——
恰是第 8 题，借助
\cref{exo:g12:complex:9}的核心来看。

\textbf{15.} $8\iu = 8\eu^{\iu\pi/2}$：根为
$2\eu^{\iu\pi/6} = \sqrt3 + \iu$；
$2\eu^{5\iu\pi/6} = -\sqrt3 + \iu$；
$2\eu^{3\iu\pi/2} = -2\iu$。

\textbf{16.} $1 + \iu = \sqrt2\,\eu^{\iu\pi/4}$：该映射
将平面旋转 $45^\circ$ 并放大 $\sqrt2$ 倍。单位
正方形变为边长为 $\sqrt2$ 的斜置正方形——
面积变为两倍，旋转了八分之一圈。

\textbf{17.} $\abs{z_1 z_2}^2 = z_1 z_2 \overline{z_1 z_2} =
z_1\overline{z_1}\, z_2\overline{z_2} =
\abs{z_1}^2\abs{z_2}^2$：两边开平方。
$1 + \iu$ 的幂：$1 + \iu$、$2\iu$、$-2 + 2\iu$、$-4$：每一次
都比前一次多转 $45^\circ$，模长变为 $\sqrt2$ 倍——
形成一条向外的螺旋线，其模依次为
$\sqrt2, 2, 2\sqrt2, 4, \dots$。

\textbf{18.} $\mathrm{j}$ 是三次\omterm{def:g12:complex:rootsofunity}{单位根}，
且异于 $1$，故（沿用第 7 题的论证，取 $n = 3$）
$1 + \mathrm{j} + \mathrm{j}^2 = 0$。标准三角形上的
判别条件为：$a + \mathrm{j}b + \mathrm{j}^2 c = 1 +
\mathrm{j}\cdot\mathrm{j} + \mathrm{j}^2 \cdot \mathrm{j}^2 =
1 + \mathrm{j}^2 + \mathrm{j}^4 = 1 + \mathrm{j}^2 +
\mathrm{j} = 0$：满足，与等边三角形
$\left(1, \mathrm{j}, \mathrm{j}^2\right)$ 的要求相吻合。

\textbf{19.} 在 $\C$ 上，实系数多项式可分解为一次因式的乘积；
其非实根成共轭对出现（对\omterm{def:g10:algebra:equation}{方程}取共轭即可见），
每对相乘得一个实系数二次式
$z^2 - 2\operatorname{Re}(z_0)\,z + \abs{z_0}^2$，
于是得到\omterm{def:g10:numbers:sets}{实数}范围内由 $1$ 次与 $2$ 次因式组成的分解——
第 12 题对 $z^4 + 4$ 作了这一分解。

\textbf{20.} 本章要点：点；旋转与伸缩；正多边形；
借\omterm{def:g10:algebra:expand}{展开}得到的三角恒等式；以及一个封闭的世界，
其中每个 $n$ 次\omterm{def:g10:algebra:equation}{方程}恰有 $n$ 个根。尾声：
第 10 题中，$\pi$ 的五边形与黄金
分割握了手——欧几里得最著名的两位来宾，
终于借五个\omterm{def:g11:vect:vector}{向量}之和为零这一和式登场。
\end{solution}
